\documentclass[10pt]{article} 
\renewcommand{\thempfootnote}{\arabic{mpfootnote}}

\usepackage[utf8]{inputenc}

\usepackage{multicol}
\setlength{\columnsep}{1cm}

\usepackage{abc}
\usepackage{graphicx}
\usepackage{musicography}

\usepackage[margin=0.5in]{geometry} 
\geometry{letterpaper}

\usepackage{fancyhdr} 
\pagestyle{fancy} % options: empty , plain , fancy
\renewcommand{\headrulewidth}{0pt} % customise the layout...
\lhead{}\chead{}\rhead{}
\lfoot{}\cfoot{\thepage}\rfoot{}

\begin{document}

\thispagestyle{empty}

\begin{abc}[name=skippincat,options={-O= +c -F single.fmt}]
X:1
T:Skippin' Cat
C:Jack Link
M:C|
Z:abc-transcription: Josh Larios <jdlarios@gmail.com>, 2021.08.31
K:G
B,C | "G"DCB,C DCB,C | DGBG "C"c2 Bc | "G"dedc BAG=F | "C"EC-C2 C2 B,C | 
"G"DCB,C DCB,C | DGBG "C"cdef | "D"gage dBAc | "G"BG-G2 G2 :|
|: z2 | "D"A2-Ac "G"B2 d2 | "D"ABcA "G"BG-G2 | "D"A2-Ac Bdef | gedc "G"BG-G2 | 
"D"A2-Ac "G"B2 d2 | "D"ABcA "G"Bdef | "D"gage dBAc |"G"BG-G2 G2 :|

%%center © Copyright 1986, Jack Link. All Rights Reserved. Used by Permission.
\end{abc}
\vskip 0.5in
\noindent\makebox[\textwidth][c]{%
  \begin{minipage}{7.0in}
   \setlength{\parindent}{1.5em}

   Skippin' Cat is a tune written by Jack Link, who is probably best
   known as a member of the Gypsy Gyppo String Band and the first
   incarnation of the Hurricane Ridgerunners, among other influential
   Seattle-area old-time bands in the 1970s and on. The Gyppos' Monday
   Night Dance is one of the seeds from which the current traditional
   square and contra dancing scene still going strong in the area
   (albeit virtually during the pandemic) grew. Jack and his friends
   were also responsible for the creation of the tutorial program at
   the Festival of American Fiddle Tunes in Port Townsend, when they
   realized that there was no program for beginners at the first
   festival.

   In the notes for this tune in The Portland
   Collection\footnote{https://www.theportlandcollection.com/}, Jack
   was quoted as saying: ``Written in 1978 as a hillbilly tune, but
   it's circulated at least as much among contra dance players. Just
   goes to show you...''

   When asked in an interview\footnote{https://youtu.be/yRbxZ30XGQs}
   what sort of legacy he would like to leave in the music community
   in the Northwest, this was his reply:
   
   ``I would like to think that after I'm gone, more people will feel
   more like they have permission to mess with the tunes. These are
   not pieces of paper with Mozart's name in the upper right corner;
   these are tunes that belong to you from the minute you get
   interested in them until the minute you die. People can take them
   apart and put them back together to suit their own playing, and I
   wish more people would do it. I'm trying to get that idea
   circulating better. Many people never quite get past just learning
   one way to play a tune. That's a good place to start, if you ask
   me. Take it from there and make the tune so that it suits your own
   way of playing and your own soul. You're not constrained to one set
   of notes in a tune. I'd like to leave a legacy of permission to
   mess with the melody.''

   So keep that in mind as you play this (or any other!) tune.

   \vskip 0.5in

   Our thanks to Jack's son, Nathan, for permission to use the tune here.
   
   
 \end{minipage}
}


\end{document}
